\chapter{算术}\label{ch:g12:arith}

算术研究整数：整除、质数、余数。
长期被视为最纯粹的数学，如今它守护着每一笔
网络支付：RSA 密码体制依赖于本章证明的贝祖、高斯与
费马定理。

\section{整除与带余除法}

\begin{definition}[整除]\label{def:g12:arith:divides}
设 $a, b \in \Z$。称 $b$ \emph{整除} $a$\index{整除}，记作
$b \mid a$，若存在 $k \in \Z$ 使 $a = kb$，也称 $a$ 为 $b$ 的
\emph{倍数}。
\end{definition}

\begin{proposition}\label{prop:g12:arith:divprops}
若 $c \mid a$ 且 $c \mid b$，则 $c$ 整除每个
形如 $au + bv$（$u, v \in \Z$）的整数组合。若 $a \mid b$ 且 $b \mid a$，其中
$a,b \in \N$，则 $a = b$。若 $a \mid b$ 且 $b \neq 0$，则 $\abs a \leq \abs b$。
\end{proposition}

\begin{proof}
写 $a = kc$，$b = lc$，则 $au + bv = (ku + lv)c$。其余结论
由 $\abs{a} = \abs{k}\,\abs{b}$ 得出，其中
$\abs k \geq 1$，当 $b = ka \neq 0$ 时。
\end{proof}

\begin{theorem}[带余除法]\label{thm:g12:arith:euclid}
\index{带余除法}
设 $a \in \Z$，$b \in \N^*$。存在唯一的数对
$(q, r) \in \Z \times \N$，使得
\[
a = bq + r \qquad\text{且}\qquad 0 \leq r < b .
\]
$q$ 称为\emph{商}，$r$ 称为\emph{余数}。
\end{theorem}

\begin{proof}
\emph{存在性。}$b$ 的倍数中不超过 $a$ 者构成的集合
有最大元 $bq$（该集合非空且有上界），令 $r = a - bq$。
由最大性，$b(q+1) > a$，从而 $0 \leq r < b$。
\emph{唯一性。}若 $bq + r = bq' + r'$，其中 $0 \leq r, r' < b$，则
$b(q - q') = r' - r$，且 $\abs{r' - r} < b$：$b$ 的倍数中绝对值
小于 $b$ 者必为 $0$，于是 $r = r'$，$q = q'$。
\end{proof}

\section{同余}

\begin{definition}[同余]\label{def:g12:arith:congruence}
设 $n \in \N^*$。称整数 $a, b$ 模
$n$ \emph{同余}\index{同余}，记作 $a \equiv b \pmod n$，若 $n \mid (a - b)$
--- 等价地，若 $a$ 与 $b$ 在
以 $n$ 作带余除法时余数相同。
\end{definition}

\begin{proposition}[运算的相容性]\label{prop:g12:arith:congops}
若 $a \equiv b \pmod n$ 且 $c \equiv d \pmod n$，则
\[
a + c \equiv b + d, \qquad
ac \equiv bd, \qquad
a^k \equiv b^k \ (k \in \N) \pmod n .
\]
\end{proposition}

\begin{proof}
$n$ 整除 $(a-b) + (c-d) = (a+c) - (b+d)$，且
$ac - bd = a(c - d) + d(a - b)$ 也是 $n$ 的倍数。幂的规则
由乘法规则用归纳法推出。
\end{proof}

\begin{method}[计算模 $n$ 的幂]\label{met:g12:arith:powers}
计算 $a^k \bmod n$ 时，先将底数模 $n$ 约化，再寻找
$a$ 的某个较小的幂与 $\pm1$ 同余，并用它化简指数。例如
求 $2^{100} \bmod 7$：由 $2^3 = 8 \equiv 1 \pmod 7$ 及
$100 = 3\times33 + 1$,
\[
2^{100} = \left(2^{3}\right)^{33} \times 2 \equiv 1^{33}\times 2 = 2 \pmod 7 .
\]
\end{method}

\section{最大公因数、贝祖与高斯}

\begin{definition}[最大公因数]\label{def:g12:arith:gcd}
设 $a, b$ 为不全为零的整数。\emph{最大
公因数}\index{最大公因数} $\gcd(a, b)$ 是同时整除 $a$
与 $b$。当 $\gcd(a,b) = 1$ 时，称 $a$ 与 $b$
\emph{互质}\index{互质}。
\end{definition}

\begin{proposition}[欧几里得算法]\label{prop:g12:arith:euclidalgo}
\index{欧几里得算法}
若 $a = bq + r$（$b \neq 0$），则 $\gcd(a, b) = \gcd(b, r)$。反复进行
带余除法，因此可求出 $\gcd(a,b)$：最大公因数即为最后一个
非零余数。
\end{proposition}

\begin{proof}
$a$ 与 $b$ 的任一公因数都整除 $r = a - bq$
（由\cref{prop:g12:arith:divprops}），从而是 $b$ 与 $r$ 的公因数；
反之，由 $a = bq + r$ 可知。两对数有相同的公因数，
故最大公因数相同。算法之所以终止，是由于余数
构成严格递减的非负整数数列。
\end{proof}

\begin{example}\label{ex:g12:arith:euclidalgo}
$\gcd(252, 198)$: $252 = 198 + 54$; $198 = 3\times54 + 36$;
$54 = 36 + 18$；$36 = 2 \times 18 + 0$。因此 $\gcd(252,198) = 18$。
\end{example}

\begin{theorem}[贝祖等式]\label{thm:g12:arith:bezout}
\index{贝祖等式}
设 $a, b$ 为不全为零的整数，$d = \gcd(a,b)$。存在
$u, v \in \Z$，使得
\[
au + bv = d .
\]
特别地，$a$ 与 $b$ 互质，当且仅当 $au + bv = 1$
对某些整数 $u, v$ 成立。
\end{theorem}

\begin{proof}
将欧几里得算法反向进行：每个余数都是前两个数的整数组合，
而初始数据 $a, b$ 也是它们自身的组合；
由逐步代入可知，最后一个非零余数 $d$ 是
$a$ 与 $b$ 的整数组合。（在\cref{ex:g12:arith:euclidalgo}中：
$18 = 54 - 36 = 54 - (198 - 3\times54) = 4\times54 - 198 =
4(252 - 198) - 198 = 4\times252 - 5\times198$。）

对于等价性：若 $\gcd(a,b) = 1$，由贝祖等式可得 $u, v$；反之，
$a$ 与 $b$ 的任一公因数都整除 $au + bv = 1$，
从而 $\gcd(a,b) = 1$。
\end{proof}

\begin{theorem}[高斯引理]\label{thm:g12:arith:gauss}
\index{高斯引理}
设 $a, b, c \in \Z$，若 $a \mid bc$ 且 $\gcd(a, b) = 1$，则 $a \mid c$。
\end{theorem}

\begin{proof}
由贝祖等式得 $au + bv = 1$，两边乘以 $c$ 得 $acu + bcv = c$。左边两项
都是 $a$ 的倍数（第二项因 $a \mid bc$），从而
$c$ 亦具有此性质。
\end{proof}

\begin{corollary}\label{cor:g12:arith:coprimeprod}
若 $a \mid c$，$b \mid c$ 且 $\gcd(a,b) = 1$，则 $ab \mid c$。
\end{corollary}

\begin{proof}
设 $c = ak$。由 $b \mid ak$ 且 $\gcd(a,b)=1$，依高斯引理
得 $b \mid k$，记 $k = bl$，于是 $c = abl$。
\end{proof}

\section{质数}

\begin{definition}[质数]\label{def:g12:arith:prime}
称整数 $p \geq 2$ 为\emph{质数}\index{质数}，若它的
正因数只有 $1$ 与 $p$。
\end{definition}

\begin{proposition}\label{prop:g12:arith:primedivides}
每个整数 $n \geq 2$ 都有一个质因数；若 $n$ 不是质数，则它有一个
质因数 $\leq \sqrt n$。若质数 $p$ 整除乘积 $ab$，则
$p \mid a$ 或 $p \mid b$（\emph{欧几里得引理}）。
\end{proposition}

\begin{proof}
$d \geq 2$ 为 $n$ 的最小因数，它是质数（$d$ 的任一真因数
都会是 $n$ 的更小因数）。若 $n = de$ 为合数，且
$2 \leq d \leq e$，则 $d^2 \leq de = n$，从而 $d \leq \sqrt n$。
对欧几里得引理：若 $p \nmid a$，则 $\gcd(p, a) = 1$（$p$ 的因数只有
$1$ 与 $p$），由高斯引理得 $p \mid b$。
\end{proof}

\begin{theorem}[欧几里得]\label{thm:g12:arith:infiniteprimes}
质数有无穷多个。
\end{theorem}

\begin{proof}
设 $p_1, \dots, p_k$ 为任意有限个质数，考虑
$N = p_1 p_2 \cdots p_k + 1$。某个质数 $p$ 整除 $N$；但每个 $p_i$
都不整除 $N$（余数为 $1$），因此 $p$ 是不在该表中的质数。
因此任何有限个质数都不能穷尽所有质数。
\end{proof}

\begin{theorem}[算术基本定理]\label{thm:g12:arith:fta}
\index{算术基本定理}
每个整数 $n \geq 2$ 都可以写成若干质数之积，且该分解
在不计因子次序的意义下唯一：
\[
n = p_1^{\alpha_1} p_2^{\alpha_2} \cdots p_r^{\alpha_r},
\qquad p_1 < p_2 < \dots < p_r \text{ 为质数},\ \alpha_i \geq 1 .
\]
\end{theorem}

\begin{proof}
\emph{存在性}，用强归纳法：若 $n$ 为质数，则它本身即为其分解；
否则 $n = de$，其中 $2 \leq d, e < n$，由归纳假设，两者都可分解
为质数之积。\emph{唯一性}：设
$p_1\cdots p_s = q_1 \cdots q_t$（均为质数，允许重复）。由
欧几里得引理，$p_1$ 整除某个 $q_j$，二者皆为质数，故 $p_1 = q_j$；
约去该因子并重复上述过程，两种分解逐项相同。
\end{proof}

\begin{theorem}[费马小定理]\label{thm:g12:arith:fermat}
\index{费马小定理}
设 $p$ 为质数，$a \in \Z$ 且 $p \nmid a$。则
\[
a^{p-1} \equiv 1 \pmod p .
\]
对每个 $a \in \Z$（无需假设互质），$a^p \equiv a \pmod p$。
\end{theorem}

\begin{proof}
考虑 $p - 1$ 个整数 $a, 2a, 3a, \dots, (p-1)a$ 模 $p$ 的情形。其中
没有一个 $\equiv 0$（若 $p \mid ka$，$1 \leq k \leq p-1$，则由欧几里得引理
得 $p \mid k$，矛盾），且它们模 $p$ 两两不同余（若
$ka \equiv la$，则 $p \mid (k - l)a$，从而 $p \mid k - l$，故 $k = l$）。
因此模 $p$ 意义下，它们恰为 $1, 2, \dots, p-1$ 的某种排列。
将所有同余式相乘得：
\[
a^{p-1}\,(p-1)! \equiv (p-1)! \pmod p .
\]
由于 $p$ 不整除 $1, \dots, p-1$ 中的任何数，反复应用欧几里得引理，
可约去 $(p-1)!$，得 $a^{p-1} \equiv 1$。第二个式子
由上式两边乘以 $a$ 即得（当 $p \mid a$ 时显然成立）。
\end{proof}

\begin{example}[密码学应用]\label{ex:g12:arith:rsa}
费马小定理使模 $n$ 的幂运算在指数选取得当时可逆——
这正是\emph{RSA}密码体制的核心。
取大质数 $p, q$，令 $n = pq$，公开 $n$ 与一个指数
$e$；加密映射为 $x \mapsto x^e \bmod n$。解密需要指数
$d$，满足 $ed \equiv 1 \pmod{(p-1)(q-1)}$，只有知道 $p$
与 $q$ 的人才能算出——而要恢复 $p, q$，就须
分解一个位数达数百位的整数 $n$，目前没有已知算法能在
可接受的时间内做到。
\end{example}

\section{练习}

\begin{exercise}[$\star$]\label{exo:g12:arith:1}
计算 $2026$ 除以 $17$ 的欧几里得除法的商与余数，
以及 $-2026$ 除以 $17$ 的商与余数。
\end{exercise}

\begin{exercise}[$\star$]\label{exo:g12:arith:2}
$7^{100}$ 模 $10$ 的余数是多少？（即
$7^{100}$ 的末位数字是多少？）
\end{exercise}

\begin{exercise}[$\star$]\label{exo:g12:arith:3}
用欧几里得算法计算 $\gcd(1071, 462)$，并求整数
$u, v$ 使 $1071u + 462v = \gcd(1071, 462)$。
\end{exercise}

\begin{exercise}[$\star$]\label{exo:g12:arith:4}
证明：对每个 $n \in \Z$，$n^2$ 同余于 $0$ 或 $1$。
（模 $4$）。由此推出，满足 $\equiv 3 \pmod 4$ 的整数绝不可能是两个平方数
之和。
\end{exercise}

\begin{exercise}[$\star\star$]\label{exo:g12:arith:5}
证明：对所有 $n \in \N$，$n(n+1)(2n+1)$ 能被 $6$ 整除。
\end{exercise}

\begin{exercise}[$\star\star$]\label{exo:g12:arith:6}
在 $\Z$ 中解同余式 $5x \equiv 3 \pmod{11}$。
（提示：求 $5$ 模 $11$ 的逆元。）
\end{exercise}

\begin{exercise}[$\star\star$]\label{exo:g12:arith:7}
在 $\Z \times \Z$ 中解丢番图方程
\[
17x - 40y = 1,
\]
再写出 $17x - 40y = 6$ 的全部解。
\end{exercise}

\begin{exercise}[$\star\star$]\label{exo:g12:arith:8}
证明 $\sqrt2$ 是无理数，利用质因数分解的唯一性
（比较 $2$ 在 $a^2 = 2b^2$ 两边的指数）。
\end{exercise}

\begin{exercise}[$\star\star\star$]\label{exo:g12:arith:9}
设 $p$ 为质数。
\begin{enumerate}
  \item 证明：对 $1 \leq k \leq p - 1$，$p$ 整除 $\dbinom{p}{k}$。
        （提示：利用 $k\binom pk = p\binom{p-1}{k-1}$、\cref{exo:g12:comb:7}
        以及高斯引理。）
  \item 对 $a \geq 0$ 用归纳法，给出费马小定理的另一证明，
        形式为 $a^p \equiv a \pmod p$。
\end{enumerate}
\end{exercise}

\begin{exercise}[$\star\star\star$]\label{exo:g12:arith:10}
\emph{（中国剩余定理问题。）}求所有使下列各式成立的整数 $n$：
\[
n \equiv 2 \pmod 3, \qquad n \equiv 3 \pmod 5, \qquad n \equiv 2 \pmod 7 .
\]
（提示：先解前两个条件，再代入第三个条件；贝祖
系数会有帮助。）
\end{exercise}

\section{问题：密码与校验码}

\begin{problem}[{周末问题 --- 同余式守护着每个条形码与每张信用卡，
而费马小定理掌管着
世界机密的锁}]
\label{pb:g12:arith:1}
G.\,H.~哈代于 1940 年曾自夸，数论“未受应用玷污”。
八十年后，每一声条形码的“嘀”响、每一笔
信用卡支付与每一条加密信息，都与他唱反调
--- 正是依靠本章的工具：同余式
（\cref{prop:g12:arith:congops}）、贝祖逆元、
（\cref{thm:g12:arith:bezout}）以及费马小定理
（\cref{exo:g12:arith:9}）。本题检验这些编码，破解一把
简化版的锁，并探究真正的锁为何安全。

\medskip
\noindent\textbf{第一部分 --- 同余运算熟练度。}
\begin{enumerate}
  \item 计算 $2026 \bmod 7$；再求
        $7^{100}$ 的末位数字（先求 $7$ 的幂模
        $10$ 的周期）。
  \item 快速幂（\cref{met:g12:arith:powers}）：
        计算 $5^{117} \bmod 13$（从
        $5^2 \equiv -1$ 出发）。
  \item 解同余式 $3x \equiv 5 \pmod 7$。
  \item 对 $(97, 35)$ 运行欧几里得算法，回代
        求整数 $u, v$ 使 $97u + 35v = 1$，
        并由此推出 $35$ 模 $97$ 的逆元。
  \item 准确叙述 $a$ 模 $n$ 可逆的条件，
        并说明由哪个定理给出其逆元。
\end{enumerate}

\medskip
\noindent\textbf{第二部分 --- 校验码。}
\begin{enumerate}[resume]
  \item ISBN-10：一本书编码的十个数字 $d_1 \dots d_{10}$
        须满足下式
        $10d_1 + 9d_2 + \dots + 2d_9 + 1d_{10} \equiv 0
        \pmod{11}$。验证真实的 ISBN
        $0\,306\,40615\,2$.
  \item 证明 ISBN 方案能检测\emph{每一个}
        单个数字错误：若某一位数字改变了 $d \not\equiv 0$，
        则加权和改变了 $w d$，
        其中 $1 \leq w \leq 10$ --- 为何它永远不可能
        $\equiv 0 \pmod{11}$
        （\cref{thm:g12:arith:gauss}）？
  \item 证明它也能检测任意两个相邻（且不同）数字的
        对调。再说明设计上的关键：
        $11$ 具有什么性质使两个证明都成立，
        若取模数 $10$，又会出什么问题？
  \item EAN-13 条形码对各位数字赋予权重 $1, 3, 1, 3, \dots$
        并取模 $10$，求使
        $978\,2940199\,05$ 完整的校验位。下列相邻对调中，
        哪些是 EAN \emph{无法}检测的？（何时
        $2(a - b) \equiv 0 \pmod{10}$？）
  \item 信用卡采用 Luhn 方案：从右往左，每隔一位数字加倍
        （减去 $9$，条件是加倍后的数超过 $9$），
        再将所有数字相加，并要求和为
        $10$ 的倍数。验证测试号码
        $4539\,1488\,0343\,6467$.
  \item 用一句话说明：采用质数模数，使 ISBN 获得了什么，
        而受限于模 $10$ 的 EAN 与 Luhn 方案所不具备的性质？
\end{enumerate}

\medskip
\noindent\textbf{第三部分 --- 费马之锁。}
\begin{enumerate}[resume]
  \item 宝藏之前的陷阱：计算
        $2^{10} \bmod 341$，由此推出 $2^{340} \bmod 341$ ---
        然后将 $341$ 作因数分解。这个例子（一个
        \emph{费马伪素数}）对于把费马
        小定理用作素性检验有何启示？
  \item 迷你 RSA：取 $p = 3$，$q = 11$，于是
        $n = 33$，$(p-1)(q-1) = 20$；公钥指数为
        $e = 3$。求私钥指数 $d$，使得
        $3d \equiv 1 \pmod{20}$（用第 4 题的方法）。
  \item 加密消息 $m = 4$：计算
        $c = m^3 \bmod 33$。
  \item 解密：计算 $c^d \bmod 33$（取
        $c \equiv -2 \pmod{33}$），并还原消息。
  \item 解密为何总是成功：证明
        $m^{21} \equiv m$ 对模 $3$ 与模 $11$ 都成立，
        （在两个模中各自用费马小定理），并推出
        对模 $33$ 成立（\cref{thm:g12:arith:gauss} 将这两个
        同余式拼合起来）。$21 = ed$ 的特殊形式
        $1 + 20k$ 是从哪里来的？
  \item 这把锁的安全性：人人都知道 $n$ 与 $e$；
        要恢复 $d$ 就需要 $(p-1)(q-1)$，进而需要
        $n$ 的因数。我们的 $33$ 一眼即可分解——为什么同一方案
        在 $n$ 为六百位数时，仍能保护全世界的银行？
        （用一句话说明相乘与分解之间的
        不对称性。）
\end{enumerate}

\medskip
\noindent\textbf{第四部分 --- 经典问题。}
\begin{enumerate}[resume]
  \item 中国古代的“孙子点兵”问题（比较
        \cref{exo:g12:arith:10}）：某队士兵人数，余数为 $2$（模 $3$），
        余数为 $3$（模 $5$）。求出所有可能的人数，
        并说明
        为什么答案在模 $15$ 下唯一。
  \item 终于来了一句话的证明：由 $10 \equiv 1 \pmod 9$，
        证明每个数都与其各位数字之和同余（模 $9$）；
        由 $10 \equiv -1 \pmod{11}$ 推出
        $11$ 的交错和规则。（初中分册
        用显式代数证明了这些结论——请欣赏
        这种简洁的写法。）
  \item 终章 --- 哈代对抗条形码：回顾本章的工具
        （同余运算、贝祖逆元、费马小定理、
        拼合互质模数），以及每个工具在本题中
        何时“豁然开朗”；
        然后给出关于
        “未受玷污”的现代判定。
\end{enumerate}
\end{problem}
